\chapter{Prompts in Grasshopper}\label{ch:ghprompts}

The speech-bubble button on Grasshopper's toolbar (right, next to the V° button), or the command
\code{miniVivisection GrasshopperPrompts}, opens a window with three tabs for the definition on the canvas:
\vn{Claude}, \vn{Python} and \vn{PowerShell}. The window stays above Grasshopper's and is minimized with it.
Closing it only hides it, so the chat, Python's namespace and PowerShell keep going.

\shot{16_grasshopper_prompts}{One message in an empty definition: two sliders, a rectangle per floor, each turned by
  Twist, lofted, and a note. Claude placed the sliders, wrote the \emph{Floor Rects} script component, added Loft,
  wired them and wrote the note; its answer says what it checked, and that \keys{Ctrl+Z} takes it all back.}{fig:ghprompts}

\section{V°Claude in a definition}\label{sec:ghclaude}

Each message names the definition, whether the solver is on, and \textbf{the components selected on the canvas}.
Select what you mean, then ask: “explain these”, “add a loft between these curves and a slider for the count”.

\subsection{What Claude uses}

\begin{description}
  \item[RhinoAI's Grasshopper tools] (RhinoAI installed and its server running, \see{sec:rhinoai}): the whole canvas
        with every object's wires and messages, the component library, placing components and sliders, wiring,
        deleting, solving; the Rhino view as an image.
  \item[miniVivisection's own tools,] for what RhinoAI doesn't do: the selection in detail; a parameter's data in full,
        each item with its facts (a point's coordinates, a curve's length and whether it is closed, a Brep's area, volume
        and whether it is solid, the bounding box); a script component's code; setting sliders (value, range), Panels
        (text, Show as), toggles, value lists, a parameter's stored values and the Rhino objects it references; writing
        a script component's code with its inputs and outputs; nicknames, enabled, preview, moving, groups; removing
        wires; \textbf{Markdown notes} next to what it built; an image of the canvas.
\end{description}

Claude checks its work, the components' messages, the data, the canvas image, before it answers. Without RhinoAI,
Claude can read and set values and write scripts and notes, but not place, wire or delete components: it says what to
do instead.

\subsection{One undo step per answer}

Whatever one answer changed, components added, deleted, set, rewired, is one step in Grasshopper's undo list, named
\emph{V°Claude: (your message)}. \keys{Ctrl+Z} takes it all back and \keys{Ctrl+Y} brings it again; \code{/undo} in the chat does
the same when the last step is Claude's. The chat says what changed (\emph{2 changed, 1 added}).
Setting: \setting{One undo step per answer}.

\subsection{An answer stays with its definition}

Claude's tools work on the definition the canvas shows. Showing another one while Claude answers stops the answer, and
the chat says so, so that nothing changes the wrong definition. What it did until then is still one undo step in its own
definition. There is one conversation per definition, continued when you come back to it.

\subsection{Commands}

\code{/explain} (the selected components, or the whole definition), \code{/fix} (the components with errors or warnings),
\code{/undo}, \code{/history}, \code{/allow on|off} (RhinoAI's commands and scripts), \code{/new}, \code{/clear}, \code{/help}.
\key{Tab} completes them.

\paragraph{Read only.} With \setting{Claude may change the definition} on \emph{Read only}, Claude reads and explains, and proposes
changes in words.

\section{V°Python in Grasshopper}

While a Grasshopper Python component is paused, input runs in its paused frame, as \vn{Python} does in the Script
Editor. Otherwise it runs in the tab's own namespace, which has these helpers
(\setting{V°Python: Grasshopper helpers}):

\begin{xltabular}{\linewidth}{@{}P{0.28\linewidth}L@{}}
\toprule
\code{ghdoc()}          & the definition, a \code{GH_Document} \\
\code{ghselected()}     & the selected objects \\
\code{ghcanvas()}       & the canvas \\
\code{ghfind('Name')}   & objects by nickname \\
\code{show()}, \code{hide()}, \code{clear()}, \code{bake()} & as in the Script Editor (\see{sec:vpython}) \\
\bottomrule
\end{xltabular}

\section{V°PowerShell in Grasshopper}

A PowerShell prompt, started in the definition's folder (Documents for an unsaved one;
\setting{V°PowerShell: start in the definition's folder}).

\section{Settings}

On the tab \emph{Grasshopper}, in \vn{Settings} and \menu{Tools\>Options\>miniVivisection}: the prompts
button, each tab on or off (\vn{Claude} also needs the \vn{AI} switch, the others \vn{Debug}), whether the window stays
with Grasshopper, what Claude may change, undo, the selection, RhinoAI's tools, how many items Claude reads per
parameter, whether it may look at the canvas, \vn{Python}'s helpers and \vn{PowerShell}'s folder (\see{set:grasshopper}).
Claude Code's path, the model, starting RhinoAI's server and \code{/allow}'s default are \vn{Claude}'s settings on the
\emph{General} tab.
